%% kapa: exam=1 round=31 year=2020 period=1 form=A subject=civil_law subject_ko=민법 seq=37 number=37
%% source: Q-Net 감정평가사 기출문제 (cst00309) artlSeq=5207219 file=extracted/1008/5207219/A형/A형/2020년도 제31회 감정평가사 1차 시험_1교시_A형.pdf
%% conversion: kordoc 4.24.51 -> tools/gen1.py (auto); review=proofread-sonnet; flags=none
\begin{question}{37}
민사유치권에 관한 설명으로 옳지 않은 것은? (다툼이 있으면 판례에 따름)
\begin{choices}
  \item 수급인은 특별한 사정이 없으면 그의 비용과 노력으로 완공한 건물에 유치권을 가지지 못한다.
  \item 물건의 소유자는 그 물건을 점유하는 제3자가 비용을 지출할 때에 점유권원이 없음을 알았거나 중대한 과실로 몰랐음을 증명하여 비용상환청구권에 기한 유치권의 주장을 배척할 수 있다.
  \item 채권과 물건 사이에 견련관계가 있더라도, 그 채무불이행으로 인한 손해배상채권과 그 물건 사이의 견련관계는 인정되지 않는다.
  \item 저당권의 실행으로 부동산에 경매개시결정의 기입등기가 이루어지기 전에 유치권을 취득한 사람은 경매절차의 매수인에게 이를 행사할 수 있다.
  \item 토지 등 그 성질상 다른 부분과 쉽게 분할할 수 있는 물건의 경우, 그 일부를 목적으로 하는 유치권이 성립할 수 있다.
\end{choices}
\end{question}
\answer{3}
